\documentclass[11pt,a4paper]{article}

% --- 核心数学与样式宏包 ---
\usepackage[utf8]{inputenc}
\usepackage[T1]{fontenc}
\usepackage{amsmath, amssymb, amsthm, bm}
\usepackage{geometry}
\geometry{left=2.2cm,right=2.2cm,top=2.5cm,bottom=2.5cm}

% --- 颜色与样式定制 (简历蓝白风) ---
\usepackage[dvipsnames]{xcolor}
\usepackage{titlesec}
\usepackage{hyperref}
\usepackage[most]{tcolorbox}

% 定义深蓝色 (简历质感)
\definecolor{DeepBlue}{RGB}{0, 51, 102} 
\definecolor{LightBlue}{RGB}{235, 245, 255}

% 标题样式定制
\titleformat{\section}{\color{DeepBlue}\normalfont\Large\bfseries}{\thesection}{1em}{}
\titleformat{\subsection}{\color{DeepBlue}\normalfont\large\bfseries}{\thesubsection}{1em}{}

\hypersetup{
    colorlinks=true,
    linkcolor=DeepBlue,
    urlcolor=DeepBlue,
}

% 定制化蓝框模块 (用于定理和公式)
\newtcolorbox{mathspiritbox}[1]{
    colback=LightBlue, 
    colframe=DeepBlue,        
    fonttitle=\bfseries,
    title=#1,
    arc=1.5mm,
    boxrule=0.8pt
}

% --- 文档信息 (精简硬核版) ---
\title{
    \vspace{-2.5cm}
    \color{DeepBlue}\Huge\textbf{Math Spirit: Theoretical Foundations} \\
    \large \textit{Algorithmic Enhancements in Multi-Agent Particle Environments}
}
\author{\textbf{Rukai Guo}}
\date{\today}

\begin{document}

\maketitle

\section{Introduction \& Motivation}
In the development of our MPE project, we encountered a critical challenge: \textbf{Sparse Rewards}. While it is tempting to manually design rewards to guide agents, this often leads to unintended "Reward Hacking" behaviors.

The core motivation for our current approach stems from a fundamental question raised during technical review: 
\begin{quote}
    \textit{"How to provide guidance to the Agent without distorting the original objective? Specifically, how to prevent agents from oscillating near a target just to exploit a distance-based reward?"}
\end{quote}
This led us to re-examine the principles of \textbf{Policy Invariance}.

\section{PBRS: Theoretical Background}
To rigorously address reward hacking, we implement \textbf{Potential-Based Reward Shaping (PBRS)}, a concept formalized by \textbf{Andrew Ng, Daishi Harada, and Stuart Russell (1999)} in their seminal paper \textit{"Policy Invariance under Reward Shaping"}.

\subsection{The Ng et al. Necessary and Sufficient Condition}
Ng et al. provided a mathematical guarantee for maintaining the optimal policy $\pi^*$. They proved that a shaping function $F(s, a, s')$ preserves the optimality of a policy if and only if it takes the following form:

\begin{mathspiritbox}{Fundamental Theorem of PBRS}
    \[ F(s, a, s') = \gamma \Phi(s') - \Phi(s) \]
    where $\Phi: S \to \mathbb{R}$ is a real-valued potential function of states, and $\gamma$ is the MDP's discount factor.
\end{mathspiritbox}

\subsection{Mathematical Proof of Policy Invariance}
By substituting $F(s, a, s')$ into the Bellman optimality equation, the relationship between the original Q-value ($Q$) and the shaped Q-value ($Q'$) can be derived:
\begin{align}
    Q'(s, a) &= \mathbb{E} [r + \gamma \Phi(s') - \Phi(s) + \gamma \max_{a'} Q'(s', a')] \\
    Q'(s, a) + \Phi(s) &= \mathbb{E} [r + \gamma (Q'(s', a') + \Phi(s'))]
\end{align}
Let $\hat{Q}(s, a) = Q'(s, a) + \Phi(s)$. The equation recovers the original Bellman form. Since $\Phi(s)$ is independent of the action $a$, it follows that:
\[ \arg\max_a Q'(s, a) = \arg\max_a \hat{Q}(s, a) = \arg\max_a Q(s, a) \]
Thus, the optimal policy $\bm{\pi^*}$ remains unchanged.

\section{Ongoing Enhancements}
Beyond PBRS, the \textit{Math Spirit} repository incorporates:
\begin{itemize}
    \item \textbf{Control Variable Methods:} Reducing gradient variance in Multi-Agent training.
    \item \textbf{Automatic Control Principles:} Leveraging stability analysis to constrain agent behaviors.
\end{itemize}

\end{document}